\begin{decisione}{D0015}{V5: autoscaling dei worker sulla coda}{Proposta}{22/10/2026}

\contesto
Il corso tratta l'elasticità e l'auto scaling come risposta ai picchi di carico. In V3 i worker sono in numero fisso: o sono pochi e il picco crea arretrato, o sono tanti e restano inattivi per la maggior parte del tempo.

\scelta
V5 è V3 con KEDA, che scala i worker tra 1 e 10 in base agli ordini in coda (50 per worker), leggendo la coda ogni 5 s. Il confronto con V3 avviene sotto W2 (picco) e W3 (sovraccarico). Oltre al tempo di smaltimento si riporta il costo in \textbf{worker-secondi}. V3 si misura anche con 10 worker fissi sotto W2: così V5 si colloca tra i due estremi di V3, sia per prestazioni sia per costo.

Alternativa scartata: autoscaling sulla CPU. I worker passano il tempo ad aspettare il database, quindi la CPU non misura il lavoro arretrato.

\conseguenze
La reazione al picco ha un ritardo (lettura della coda, autoscaler, avvio dei pod) che si prevede e si misura. Il reset deve riportare i worker al minimo prima di ogni run. In locale più worker producono più capacità solo perché il tempo di servizio è simulato come attesa e non come calcolo: è una minaccia alla validità esterna, dichiarata. La campagna cresce di 30 run.

\end{decisione}
